
\subsection{以改動估計機制換取穩定：對數轉換、零格替代與有限樣本偏誤的修正}\label{sec:logzero}\label{sec:biasred}

第二類應對不補進資料，而是改動估計的機制。改動的位置有兩處：
一是把估計所在的尺度換掉，二是把估計量本身改掉。
本小節依序回顧兩者，蓋因本文的推導正落在這兩者之間。

死亡率文獻的作法屬前者，亦即\textbf{換掉對數尺度}。
\citet{brouhns2002} 以 Poisson 對數雙線性模型取代 Lee-Carter 原始的對數尺度
最小平方，\citet{delwarde2007} 在 Poisson 概似上加懲罰以平滑 $\beta_x$，
\citet{djeundje2010} 進一步放寬為容許過度離散的擬似概似，
\citet{currie2013} 則把這類受限的廣義線性模型寫成一般的計算框架。
這條路之所以可行，是因為計數尺度上零並非奇點，故零死亡格不需要任何替代值。

留在對數尺度上所付的代價，則在統計學與計量經濟學中有其獨立且完整的發展，
而該發展與死亡率文獻幾乎互不引用。以某個正值替代零、再取對數，
最早出現在比率的對數估計。\citet{haldane1956} 指出以對數勝算比估計時，
若某格為零則估計量不存在，提議在每一格加上 $1/2$；
\citet{anscombe1956} 在同一年由變異數穩定化的角度得到類似的建議。
兩者的動機都是\textbf{使估計量存在}，而非減少偏誤。
\citet{gart1967} 隨後系統比較了各種加項對偏誤與變異數的影響，
指出加項的大小會改變估計量本身，而不只是處理零的技術細節；
\citet{gart1967} 因而是文獻中第一次明確指出替代值本身即影響估計結果的研究，
但該文的設定是二項比例，而非計數資料取對數後再作雙線性配適。

轉換本身所造成的偏誤則在迴歸的脈絡中被處理。
\citet{duan1983} 指出以 $\log y$ 為反應變數配適後，把配適值取指數回推並不會
得到 $y$ 的條件期望，差額取決於誤差的分布，並提出不需指定分布的 smearing
估計量加以修正。\citet{santossilva2006} 在國際貿易的重力模型中把這一點推到
更強的結論：在異質變異之下，對數線性模型的係數估計本身即不一致，
建議改以 Poisson 擬似最大概似直接在原尺度上配適。
這兩篇的共同訊息是：取對數之後的估計量，其偏誤由誤差的分布決定，
而不是一個可以忽略的高階項。

這條線最近的發展把問題的嚴重性提到更高一層。
\citet{chenroth2024} 處理的是反應變數可能為零時所採的各種「類對數」轉換，
例如 $\log(1+y)$ 與反雙曲正弦。該文證明，此類轉換下的平均處理效果並不能
解讀為近似的百分比效果，蓋因它依賴 $y$ 的\textbf{單位}：
只要在取轉換之前先把 $y$ 的單位重新尺度化，就可以得到任意量級的效果估計。
該文進一步給出一個三難結果：當反應變數可以為零時，
不存在一個同時是個體效果之平均、單位不變、且可點識別的參數。
其所建議的替代作法之一，正是回到水準尺度並以 Poisson 迴歸表達百分比效果。
\citet{mullahy2024} 由另一個角度得到相容的結論：此類轉換把零與正值分開，
使所估得的參數與一個經過尺度調整的線性機率模型相關，
而回推後的邊際效果與彈性對轉換中的形狀參數敏感，
其量級介於未轉換的最小平方與該線性機率模型之間；
該文因而建議不要轉換反應變數，改採兩部分模型或未轉換的迴歸。
\citet{chenroth2024} 所指出的單位依賴，在 Lee-Carter 的設定下所對應的即是
零格替代值的選擇：替代值與死亡數的尺度是相對的，
改變它等價於在取對數之前改變反應變數的單位，
故該文所述「可以得到任意量級的效果」在小人口死亡率上即表現為
年齡水準參數的偏誤隨替代值而變。
值得注意的是，這個結論與死亡率文獻「換掉對數尺度」的選擇方向一致，
兩個文獻雖互不引用，卻各自得出離開對數尺度的建議。

改動的第二個位置是\textbf{把估計量本身改掉}。就小樣本估計而言，
把估計量的期望拉回真值的作法分為兩種：先取得估計值再扣除其偏誤的事後修正，
以及修改估計方程使估計值自身即為降偏誤的預先修正；
兩者皆始於概似理論內部，以下依其發展順序回顧。
\citet{bartlett1937} 先處理概似比統計量的分布，指出其與漸近卡方分布的
$O(n^{-1})$ 落差可由一個乘數修正吸收；\citet{bartlett1953} 進一步給出最大
概似估計量 $O(n^{-1})$ 偏誤的展開。\citet{coxsnell1968} 把它寫成一般的形式，
此即事後修正：先取得最大概似估計值，再扣除其偏誤的估計。
這條線的應用隨後展開，\citet{cordeiro1991} 把該修正寫成廣義線性模型中
可由一次加權最小平方求得的形式，使其在實務上可行；
該文並指出對邏吉斯迴歸而言，偏誤向量幾乎與係數向量平行，
故修正在方向上近似一個收縮。

問題則由 \citet{firth1993} 指出：事後修正必須先有有限的最大概似估計值，
而 \citet{albert1984} 早已刻畫了完全分離使該估計值發散的情形，
在這種情形下事後修正無從施加。小人口的死亡資料恰好容易落入該情形，
蓋因某個年齡若各年皆無死亡，Poisson 對數概似對該年齡的最大值即在負無窮處。
\citet{firth1993} 因而提出預先修正，把偏誤的首項移到分數函數上，
等價於在概似上乘一個 Jeffreys 型的懲罰，如此估計值恆為有限
（該文的推導另有勘誤，見 \citealp{firth1995}）。
該作法所修正的是 $O(n^{-1})$ 的首階偏誤，故修正後的估計量仍帶有高階的偏誤，
而非期望恰等於真值。

此一作法後續被廣泛接續使用，其推廣的方向涵蓋了指數族的非線性模型。
\citet{kosmidis2009} 把它擴及指數族非線性模型，使其適用範圍涵蓋 Lee-Carter
這類雙線性的設定；\citet{kosmidis2011} 以 Poisson 對數線性模型的等價表示
處理多項邏吉斯迴歸；\citet{kosmidis2014} 處理累積連結模型。
其中 \citet{kosmidis2010} 的結果最為一般：該文給出一個通用的疊代演算法，
只要某個模型的最大概似估計量的首階偏誤已經導出，即可據以計算其降偏誤估計，
並明言該演算法可視為\textbf{一連串的疊代偏誤修正}。
該文的演算法有兩項性質值得指出。首先，它把事後修正與預先修正統一為同一個
演算法的兩個端點，故「算出偏誤再扣除、並反覆至收斂」這個作法在方法論上
已有先例。其次，它所要求的輸入是「偏誤的可寫出形式」而非「概似的存在」，
故偏誤只要能由反應變數的邊際分布算得，該演算法即可施加。

後續發現的新問題有兩項。期望不偏並不蘊含中位數不偏，且前者不具再參數化
不變性：一個在某個參數化下期望不偏的估計量，經非線性重新參數化後即不再如此。
\citet{kennepagui2017} 因而提出中位數偏誤修正，同樣以修改分數方程的方式
達成，所得估計量在保持興趣參數的再參數化下不變、且為三階中位數不偏，
並與 \citet{firth1993} 同樣可避免無限估計值；
\citet{kennepagui2020} 把該作法用於 Cox 迴歸的單調概似情形。
修正的量另隨模型而異，目前並無涵蓋所有模型的通用公式；
\citet{cordeiro1991} 在結語中即以此說明每一類模型都須個別推導的原因，
而這條線至今仍在逐一擴張，例如 \citet{sossa2026} 把 Cox--Snell 修正推到
均值與離散度同時變動的異質變異非線性模型。

\citet{breslow1995} 的結果與本問題的情形最為接近。該文處理帶單一變異成分的
廣義線性混合模型，推導其三種近似估計量（含懲罰擬似概似）的漸近偏誤，
發現一階估計量對變異成分的估計有嚴重偏誤，
而該偏誤的量\textbf{由變異成分的大小決定，而不是由樣本數決定}；
該文並給出易算的修正因子。\citet{linbreslow1996} 把結果推廣到多個變異成分。
此一分支確立了一個先例：偏誤修正的原理不限於 $O(n^{-1})$ 的漸近項，
而可施加於任何可寫出其偏誤的近似，且該偏誤不必隨樣本數消失。

這一分支同時把本小節的兩處改動接了起來。\citet{lin1999} 以平滑樣條處理廣義
加法混合模型的推論，\citet{kauermann2009} 給出廣義懲罰樣條平滑的漸近結果，
\citet{xiao2020} 則處理函數型資料的懲罰樣條；
而本小節開頭所述死亡率修勻的一路
（\citealp{currie2004}；\citealp{delwarde2007}；\citealp{currie2013}）
在統計理論上即屬這一類模型。亦即死亡率文獻為了平滑而引入的工具，
其自身的文獻早已記載其近似概似會產生不隨樣本數消失的偏誤；
死亡率文獻引入這些工具時關注的是平滑度，並未援引這一側的偏誤結果。

這條脈絡所記載的結果都建立在「有一個可寫下的概似」之上，
而標準 Lee-Carter 先取對數再作最小平方，其目標函數本就不是任何分配的對數
概似，故其結果無法直接施加於它。缺口因而相當具體：
標準 Lee-Carter 的實作至今仍以對數尺度的奇異值分解為主，
\citet{lee1992} 的原始作法如此，\citet{yue2019} 在小人口的研究中所修改的
對象亦是該作法；而只要使用這個作法，零格就必須以某個正值替代。
就本文所查，該替代值的選擇對 Lee-Carter 參數估計的影響尚未被量化，
$0.5$ 這個慣例在死亡率文獻中亦未見其出處的討論；
\citet{haldane1956} 以來的「加一個常數」、\citet{duan1983} 以來的
「轉換造成偏誤」與 \citet{chenroth2024} 的「單位依賴」三條線索，
至今未被接到小人口死亡率的設定上。
